%% =====================================================================
%%  4e de couverture — Python, volume 1 (toujours sur une page paire)
%% =====================================================================
\clearpage\ifodd\value{page}\otPageNotes\clearpage\fi
\thispagestyle{empty}
\begin{tikzpicture}[remember picture,overlay]
\begin{scope}[shift={(current page.south west)}]
  \fill[otNavy] (0,0) rectangle (210mm,297mm);
  \draw[white,opacity=0.05,line width=0.3pt,step=7mm] (0,0) grid (210mm,297mm);
  \fill[otGreen] (0,0) rectangle (210mm,2.4mm);

  \node[anchor=north west,inner sep=0pt,text width=165mm,text=white,font=\otBold\fontsize{22}{27}\selectfont]
    at (24mm,262mm) {Never written a line of code?\\This book is for you.};
  \fill[otGreen] (24mm,236mm) rectangle (54mm,237.4mm);
  \node[anchor=north west,inner sep=0pt,text width=152mm,text=white,opacity=0.9,
    font=\otRegular\fontsize{11}{16}\selectfont] at (24mm,228mm)
    {Python is one of the most popular programming languages in the world: it runs websites, analyses
     data, automates boring tasks and powers artificial intelligence. This method takes you from
     your very first instruction to complete programs, one small step at a time. Every idea is
     explained in plain English, every program has been run and tested, every exercise comes with
     a detailed solution.};

  % Trois arguments
  \foreach \n/\ico/\titre/\txt in {0/livre/A complete method/20 chapters and 5 guided projects: lessons\, examples\, exercises and reviews.,
      1/clavier/Learn by doing/Hundreds of short programs to type\, run and change\, all of them downloadable.,
      2/bug/Made for beginners/No jargon left unexplained\, and the most common mistakes anticipated and fixed.}{%
    \pgfmathsetmacro{\x}{24+\n*55}
    \fill[white,opacity=0.07,rounded corners=2mm] (\x mm,140mm) rectangle ++(51mm,48mm);
    \fill[otGreen,rounded corners=1.6mm] (\x mm+4mm,174mm) rectangle ++(10mm,10mm);
    \otDessinIcone{\ico}{white}{otGreen}{6.4mm}{(\x mm+9mm,179mm)}
    \node[anchor=north west,inner sep=0pt,text width=44mm,text=white,align=left] at (\x mm+4mm,170mm)
      {{\otSemi\fontsize{10.4}{12.4}\selectfont\titre}\par\vspace{1.4mm}
       {\otRegular\fontsize{8.8}{11.4}\selectfont\color{white!80!otNavy}\txt}};}

  % Au programme
  \node[anchor=north west,inner sep=0pt,text width=162mm,text=white,align=left] at (24mm,128mm)
    {{\otSemi\fontsize{7.5}{9}\selectfont\color{otGreen}\otCaps[12]{What you will learn}}\par\vspace{2.4mm}
     {\otRegular\fontsize{10}{14}\selectfont Install Python\enspace·\enspace Calculate\enspace·\enspace
      Variables\enspace·\enspace Text and f-strings\enspace·\enspace Ask the user\enspace·\enspace
      Make decisions\enspace·\enspace Loops\enspace·\enspace Games with randomness\enspace·\enspace
      Functions\enspace·\enspace Handle errors\enspace·\enspace Lists\enspace·\enspace
      Dictionaries\enspace·\enspace Files}};

  % Échelle des niveaux
  \foreach \n/\niv in {0/1,1/2,2/3,3/4,4/5}{%
    \pgfmathsetmacro{\x}{24+\n*20}
    \ifnum\n=0 \fill[otGreen,rounded corners=1mm] (\x mm,72mm) rectangle ++(18mm,9mm);
    \else \fill[white,opacity=0.12,rounded corners=1mm] (\x mm,72mm) rectangle ++(18mm,9mm);\fi
    \node[text=white,font=\otBold\fontsize{10}{11}\selectfont] at (\x mm+9mm,76.5mm) {\niv};}
  \foreach \n/\nom in {0/Beginner,1/Elementary,2/Intermediate,3/Advanced,4/Expert}{%
    \pgfmathsetmacro{\x}{24+\n*20}
    \node[anchor=north,text=white,opacity=0.7,font=\otRegular\fontsize{6.2}{8}\selectfont] at (\x mm+9mm,70mm) {\nom};}
  \node[anchor=west,text=white,opacity=0.85,font=\otRegular\fontsize{9}{11}\selectfont,text width=60mm] at (128mm,76mm)
    {A collection in five volumes, from your first program to professional development.};

  % Pied : logo, site, emplacement du code-barres
  \node[anchor=south west,inner sep=0pt] at (24mm,24mm) {\includegraphics[width=44mm]{logo-blanc}};
  \node[anchor=base west,inner sep=0pt,text=white,opacity=0.75,font=\otMedium\fontsize{9.5}{11}\selectfont]
    at (24mm,16mm) {\otSite};
  \fill[white] (146mm,16mm) rectangle (186mm,40mm);
  \node[text=otMuted,font=\otRegular\fontsize{7}{9}\selectfont,align=center] at (166mm,28mm) {ISBN\\forthcoming};
\end{scope}
\end{tikzpicture}
